\section{方法对比与性能分析}

\subsection{MDLM vs BD3LM vs SoLM}

BD3LM（Block Diffusion Language Model）是另一种支持 KV 缓存的扩散语言模型方法。其思路是将长序列分成固定大小的块（例如每块 100 个 token），在每个块内做扩散去噪，块间则使用自回归方式推进。

\begin{table}[H]
\centering
\caption{三种扩散语言模型方法对比}
\begin{tabular}{@{}lccc@{}}
\toprule
\textbf{特性} & \textbf{MDLM} & \textbf{BD3LM} & \textbf{SoLM} \\
\midrule
注意力类型 & 双向 & 块内双向+块间因果 & 因果（排序后） \\
KV 缓存 & 不支持 & 部分支持（仅跨块） & 完全支持 \\
并行生成 & 全局并行 & 块内并行 & 全局并行 \\
推理速度 & 慢 & 中等 & 快 \\
生成质量 & 高 & 中-高 & 高 \\
\bottomrule
\end{tabular}
\end{table}

\begin{warningbox}{BD3LM 在少步采样时的质量退化}
BD3LM 在增加采样步数时能维持较好的质量，但当大幅减少采样步数（即每步并行生成更多词）时，质量会显著退化，甚至可能不如 MDLM。这是因为块内扩散的步数减少后，去噪效果不足。相比之下，SoLM 在少步采样时仍能保持较好的质量。
\end{warningbox}

\subsection{速度-质量权衡曲线}

在困惑度 vs 采样时间的 Pareto 前沿上：

\begin{itemize}
    \item \textbf{MDLM}：困惑度最优，但推理速度最慢（曲线偏右下方）
    \item \textbf{BD3LM}：在高采样步数时快于 MDLM，但减少步数后质量急剧下降
    \item \textbf{SoLM}：Pareto 前沿\textbf{严格优于} MDLM 和 BD3LM，即在任意给定速度下质量更高，或在任意给定质量下速度更快
\end{itemize}

在 110M 参数、1024 token 上下文长度的基准测试中，SoLM 相比 MDLM 的采样时间实现了显著的减少。

\subsection{本章小结}

SoLM 在速度-质量权衡上严格优于 MDLM 和 BD3LM，这主要归功于其完全的 KV 缓存支持。BD3LM 的块式设计虽然也能部分支持 KV 缓存，但块内仍受限于双向注意力的瓶颈。

\newpage
